




\begin{table*}[t]
    \centering
    \caption{
    \textbf{Ablation study on key components of \ourmodel.}
    Results are evaluated on 1.7B student model.}

    \vspace{1 mm}
    \label{tab:ablation}
    
    \begingroup
    \setlength{\tabcolsep}{4pt}
    \renewcommand{\arraystretch}{1.08}
    \normalsize

    \begin{adjustbox}{max width=0.90\textwidth}
    \begin{tabular}{lccccc}
        \toprule
        \multirow{2}{*}{\textbf{Variants}}
        & \multicolumn{2}{c}{\textbf{Math}}
        & \textbf{Science}
        & \textbf{Code}
        & \multirow{2}{*}{\textbf{Average}} \\
        \cmidrule(lr){2-3}
        \cmidrule(lr){4-4}
        \cmidrule(lr){5-5}
        & AIME 2024 & AIME 2025 & GPQA & LiveCodeBench & \\

        \midrule
        \rowcolor[HTML]{FEE090}
        \multicolumn{6}{c}{\textbf{Step-wise Reweighting}}\\
        (1.1) \textit{w/} Uniform Weighting
        & 41.68\phantom{.}\scalebox{0.72}{$\pm$0.95} & 35.58\phantom{.}\scalebox{0.72}{$\pm$1.01} & 30.12\phantom{.}\scalebox{0.72}{$\pm$0.59} & 31.43\phantom{.}\scalebox{0.72}{$\pm$0.68} & 34.70 \\
        (1.2) \textit{w/} Heuristic Weighting
        & 44.96\phantom{.}\scalebox{0.72}{$\pm$0.98} & 38.75\phantom{.}\scalebox{0.72}{$\pm$1.04} & 31.14\phantom{.}\scalebox{0.72}{$\pm$0.64} & 33.71\phantom{.}\scalebox{0.72}{$\pm$0.72} & 37.14 \\
        (1.3) \textit{w/} Mask After Wrong
        & 39.11\phantom{.}\scalebox{0.72}{$\pm$1.08} & 31.59\phantom{.}\scalebox{0.72}{$\pm$1.14} & 26.56\phantom{.}\scalebox{0.72}{$\pm$0.67} & 30.12\phantom{.}\scalebox{0.72}{$\pm$0.80} & 31.85 \\
        (1.4) \textit{w/o} Weight Clipping
        & 45.78\phantom{.}\scalebox{0.72}{$\pm$1.16} & 37.93\phantom{.}\scalebox{0.72}{$\pm$1.21} & 33.57\phantom{.}\scalebox{0.72}{$\pm$0.74} & 35.12\phantom{.}\scalebox{0.72}{$\pm$0.79} & 38.10 \\

        \midrule
        \rowcolor[HTML]{E0F3F8}
        \multicolumn{6}{c}{\textbf{Training Objective}}\\
        (2.1) \textit{w/o} GRPO
        & 48.87\phantom{.}\scalebox{0.72}{$\pm$0.93} & 39.73\phantom{.}\scalebox{0.72}{$\pm$1.02} & 35.89\phantom{.}\scalebox{0.72}{$\pm$0.62} & 38.62\phantom{.}\scalebox{0.72}{$\pm$0.71} & 40.78 \\
        (2.2) \textit{w/o} Step-wise OPD
        & 25.63\phantom{.}\scalebox{0.72}{$\pm$1.03} & 21.67\phantom{.}\scalebox{0.72}{$\pm$1.12} & 33.55\phantom{.}\scalebox{0.72}{$\pm$0.76} & 20.70\phantom{.}\scalebox{0.72}{$\pm$0.87} & 25.39 \\

        \midrule

        \textbf{\ourmodel}
        & \textbf{50.83}\phantom{.}\scalebox{0.72}{$\pm$1.15} & \textbf{41.72}\phantom{.}\scalebox{0.72}{$\pm$1.24} & \textbf{38.72}\phantom{.}\scalebox{0.72}{$\pm$0.73}
        & \textbf{40.63}\phantom{.}\scalebox{0.72}{$\pm$0.91} & \textbf{42.98} \\

        \bottomrule
    \end{tabular}%
    \end{adjustbox}

    \endgroup

    \vspace{2pt}
    \vspace{-10pt}
\end{table*}




\subsection{Ablation Study}



To evaluate each component in \ourmodel, we develop six variants from two perspectives on the Qwen3-1.7B student. For \textit{Step-wise Reweighting}, we consider: (1.1)~\textit{w/ Uniform Weighting}, which fixes $w_k=1$ for all steps; (1.2)~\textit{w/ Heuristic Weighting}, which applies a fixed exponential decay $w_k = \gamma^{k - k_{\mathrm{err}}}$ starting from the first erroneous tool call at step $k_{\mathrm{err}}$, with $\gamma$ fixed as 0.9; (1.3)~\textit{w/ Mask After Wrong}, which zeros out the OPD signal for all steps after the first tool-call error; and (1.4)~\textit{w/o Weight Clipping}, which removes the upper-bound clipping $\delta$ in~\cref{eq:adaptive_weight}. For \textit{Training Objective}, we ablate the two terms in \cref{eq:training_object}: (2.1)~\textit{w/o GRPO}, which removes $\mathcal{L}_{\mathrm{GRPO}}$; and (2.2)~\textit{w/o Step-wise OPD}, which removes $\mathcal{L}_{\mathrm{OPD}}^{\mathrm{step}}$. The results are reported in Table~\ref{tab:ablation}, from which we observe:



\begin{itemize}[leftmargin=*]

\item \textbf{Obs 3: The adaptive step-wise reweighting is critical, and neither static nor heuristic alternatives can substitute it.}
Replacing our adaptive reweighting mechanism with uniform weighting (variant 1.1) causes a notable drop to 34.70\% average, confirming that treating all steps equally exposes training to misleading supervision from corrupted states. Heuristic decay (variant 1.2) partially mitigates this issue (37.14\%) by down-weighting later steps, yet still falls behind because a fixed schedule cannot capture non-monotonic divergence patterns where the student recovers alignment after earlier errors. Hard masking (variant 1.3) performs worst (31.85\%), as it discards all subsequent supervision after a single error, wasting informative signals from partially correct trajectories and preventing the student from recovering. Removing weight clipping (variant 1.4) yields 38.10\%, indicating that unbounded amplification introduces training instability. 

\item \textbf{Obs 4: Both the GRPO and step-wise OPD components are indispensable, serving complementary roles.}
Removing GRPO (variant 2.1) reduces the average to 40.78\%, showing that outcome-level rewards remain valuable for steering exploration with dense distillation. Removing step-wise OPD (variant 2.2) causes a far more severe drop to 25.39\%, 
which confirms that sparse rewards alone are insufficient for stable TIR learning in small models. The pronounced asymmetry reveals their complementary nature: step-wise OPD provides the primary fine-grained guidance, while GRPO broadens trajectory-level exploration beyond the teacher distribution.


\end{itemize}



